% Sección para el informe de la Entrega 2. Se incluye con % Sección para el informe de la Entrega 2. Se incluye con % Sección para el informe de la Entrega 2. Se incluye con % Sección para el informe de la Entrega 2. Se incluye con \input{fase2_seccion_integracion_espacial}
% (requiere booktabs y listings, ya cargados en main.tex). Reemplace los campos <...> con la salida real de
% src/spatial_evidence.py (docs/evidencia/spatial_evidence.md).
\section{Integración Espacial en MongoDB}

\subsection{Carga e índices}
Los municipios PDET limpios (GeoJSON en EPSG:4326, producido por el ETL de integridad) se cargan en la colección
\texttt{municipalities} con \texttt{src/load\_to\_mongo.py}. La carga es idempotente (\emph{upsert} por
\texttt{\_id}, el código DIVIPOLA), las geometrías se normalizan a \texttt{MultiPolygon} con la regla de la mano
derecha y se rechazan las inválidas o ubicadas fuera de la caja de Colombia. Los índices se crean antes de insertar,
de modo que MongoDB también rechaza cualquier geometría que no pueda indexar.

\begin{lstlisting}[style=lab,language=mongosh]
db.municipalities.createIndex({ geometry: "2dsphere" }, { name: "geometry_2dsphere" })
db.municipalities.createIndex({ pdet_subregion: 1 },    { name: "pdet_subregion_1" })
\end{lstlisting}

La carga queda registrada en \texttt{ingestion\_log} (archivo, filas leídas, insertadas y rechazadas, SHA-256 del
archivo y fecha). Resultado: <N> municipios cargados de <N> esperados (170), <S> subregiones (16), <R> rechazados.

\subsection{Evidencia del uso de los índices}
Se ejecutaron consultas \texttt{\$geoIntersects}, \texttt{\$geoWithin} y \texttt{\$near} con
\texttt{explain("executionStats")}, comparando el plan con el índice \texttt{2dsphere} frente a un recorrido
completo forzado con \texttt{hint(\{\$natural: 1\})}. El Cuadro~\ref{tab:evidencia} resume los planes ganadores y
los documentos examinados.

\begin{table}[!htbp]
\centering\small
\input{../evidencia/spatial_evidence_table}   % generado por src/spatial_evidence.py
\caption{Planes de ejecución de las consultas espaciales sobre \texttt{municipalities}.}
\label{tab:evidencia}
\end{table}

Los controles automáticos (índice presente, número de municipios y subregiones, códigos DIVIPOLA de cinco
dígitos, tipo \texttt{MultiPolygon}, y consultas punto-en-polígono cuyo resultado se verifica contra el municipio
de origen) se reportan en \texttt{docs/evidencia/spatial\_evidence.md}. <Comentar aquí: con 170 documentos la
diferencia de tiempo es pequeña, pero el plan cambia de \texttt{COLLSCAN} a \texttt{IXSCAN}; la ventaja crece en
las colecciones de edificios (Entrega 3), con millones de documentos.>

% (requiere booktabs y listings, ya cargados en main.tex). Reemplace los campos <...> con la salida real de
% src/spatial_evidence.py (docs/evidencia/spatial_evidence.md).
\section{Integración Espacial en MongoDB}

\subsection{Carga e índices}
Los municipios PDET limpios (GeoJSON en EPSG:4326, producido por el ETL de integridad) se cargan en la colección
\texttt{municipalities} con \texttt{src/load\_to\_mongo.py}. La carga es idempotente (\emph{upsert} por
\texttt{\_id}, el código DIVIPOLA), las geometrías se normalizan a \texttt{MultiPolygon} con la regla de la mano
derecha y se rechazan las inválidas o ubicadas fuera de la caja de Colombia. Los índices se crean antes de insertar,
de modo que MongoDB también rechaza cualquier geometría que no pueda indexar.

\begin{lstlisting}[style=lab,language=mongosh]
db.municipalities.createIndex({ geometry: "2dsphere" }, { name: "geometry_2dsphere" })
db.municipalities.createIndex({ pdet_subregion: 1 },    { name: "pdet_subregion_1" })
\end{lstlisting}

La carga queda registrada en \texttt{ingestion\_log} (archivo, filas leídas, insertadas y rechazadas, SHA-256 del
archivo y fecha). Resultado: <N> municipios cargados de <N> esperados (170), <S> subregiones (16), <R> rechazados.

\subsection{Evidencia del uso de los índices}
Se ejecutaron consultas \texttt{\$geoIntersects}, \texttt{\$geoWithin} y \texttt{\$near} con
\texttt{explain("executionStats")}, comparando el plan con el índice \texttt{2dsphere} frente a un recorrido
completo forzado con \texttt{hint(\{\$natural: 1\})}. El Cuadro~\ref{tab:evidencia} resume los planes ganadores y
los documentos examinados.

\begin{table}[!htbp]
\centering\small
\begin{tabular}{p{0.34\linewidth}p{0.28\linewidth}rrr}
\toprule
Consulta & Plan ganador & Docs (índice) & Docs (sin índice) & Devueltos \\
\midrule
\texttt{punto\_en\_poligono[05031]} & FETCH/IXSCAN & 1 & 170 & 1 \\
\texttt{punto\_en\_poligono[13688]} & FETCH/IXSCAN & 1 & 170 & 1 \\
\texttt{punto\_en\_poligono[19548]} & FETCH/IXSCAN & 2 & 170 & 1 \\
\texttt{punto\_en\_poligono[47001]} & FETCH/IXSCAN & 1 & 170 & 1 \\
\texttt{punto\_en\_poligono[54800]} & FETCH/IXSCAN & 1 & 170 & 1 \\
\texttt{punto\_en\_el\_oceano} & FETCH/IXSCAN & 0 & 170 & 0 \\
\texttt{ventana\_geoWithin} & FETCH/IXSCAN & 34 & 170 & 18 \\
\texttt{vecinos\_geoIntersects} & FETCH/IXSCAN & 5 & 170 & 1 \\
\texttt{cercanos\_near} & GEO\_NEAR\_2DSPHERE/FETCH/IXSCAN/FETCH/IXSCAN/FETCH/IXSCAN & 335 & -- & 18 \\
\bottomrule
\end{tabular}
   % generado por src/spatial_evidence.py
\caption{Planes de ejecución de las consultas espaciales sobre \texttt{municipalities}.}
\label{tab:evidencia}
\end{table}

Los controles automáticos (índice presente, número de municipios y subregiones, códigos DIVIPOLA de cinco
dígitos, tipo \texttt{MultiPolygon}, y consultas punto-en-polígono cuyo resultado se verifica contra el municipio
de origen) se reportan en \texttt{docs/evidencia/spatial\_evidence.md}. <Comentar aquí: con 170 documentos la
diferencia de tiempo es pequeña, pero el plan cambia de \texttt{COLLSCAN} a \texttt{IXSCAN}; la ventaja crece en
las colecciones de edificios (Entrega 3), con millones de documentos.>

% (requiere booktabs y listings, ya cargados en main.tex). Reemplace los campos <...> con la salida real de
% src/spatial_evidence.py (docs/evidencia/spatial_evidence.md).
\section{Integración Espacial en MongoDB}

\subsection{Carga e índices}
Los municipios PDET limpios (GeoJSON en EPSG:4326, producido por el ETL de integridad) se cargan en la colección
\texttt{municipalities} con \texttt{src/load\_to\_mongo.py}. La carga es idempotente (\emph{upsert} por
\texttt{\_id}, el código DIVIPOLA), las geometrías se normalizan a \texttt{MultiPolygon} con la regla de la mano
derecha y se rechazan las inválidas o ubicadas fuera de la caja de Colombia. Los índices se crean antes de insertar,
de modo que MongoDB también rechaza cualquier geometría que no pueda indexar.

\begin{lstlisting}[style=lab,language=mongosh]
db.municipalities.createIndex({ geometry: "2dsphere" }, { name: "geometry_2dsphere" })
db.municipalities.createIndex({ pdet_subregion: 1 },    { name: "pdet_subregion_1" })
\end{lstlisting}

La carga queda registrada en \texttt{ingestion\_log} (archivo, filas leídas, insertadas y rechazadas, SHA-256 del
archivo y fecha). Resultado: <N> municipios cargados de <N> esperados (170), <S> subregiones (16), <R> rechazados.

\subsection{Evidencia del uso de los índices}
Se ejecutaron consultas \texttt{\$geoIntersects}, \texttt{\$geoWithin} y \texttt{\$near} con
\texttt{explain("executionStats")}, comparando el plan con el índice \texttt{2dsphere} frente a un recorrido
completo forzado con \texttt{hint(\{\$natural: 1\})}. El Cuadro~\ref{tab:evidencia} resume los planes ganadores y
los documentos examinados.

\begin{table}[!htbp]
\centering\small
\begin{tabular}{p{0.34\linewidth}p{0.28\linewidth}rrr}
\toprule
Consulta & Plan ganador & Docs (índice) & Docs (sin índice) & Devueltos \\
\midrule
\texttt{punto\_en\_poligono[05031]} & FETCH/IXSCAN & 1 & 170 & 1 \\
\texttt{punto\_en\_poligono[13688]} & FETCH/IXSCAN & 1 & 170 & 1 \\
\texttt{punto\_en\_poligono[19548]} & FETCH/IXSCAN & 2 & 170 & 1 \\
\texttt{punto\_en\_poligono[47001]} & FETCH/IXSCAN & 1 & 170 & 1 \\
\texttt{punto\_en\_poligono[54800]} & FETCH/IXSCAN & 1 & 170 & 1 \\
\texttt{punto\_en\_el\_oceano} & FETCH/IXSCAN & 0 & 170 & 0 \\
\texttt{ventana\_geoWithin} & FETCH/IXSCAN & 34 & 170 & 18 \\
\texttt{vecinos\_geoIntersects} & FETCH/IXSCAN & 5 & 170 & 1 \\
\texttt{cercanos\_near} & GEO\_NEAR\_2DSPHERE/FETCH/IXSCAN/FETCH/IXSCAN/FETCH/IXSCAN & 335 & -- & 18 \\
\bottomrule
\end{tabular}
   % generado por src/spatial_evidence.py
\caption{Planes de ejecución de las consultas espaciales sobre \texttt{municipalities}.}
\label{tab:evidencia}
\end{table}

Los controles automáticos (índice presente, número de municipios y subregiones, códigos DIVIPOLA de cinco
dígitos, tipo \texttt{MultiPolygon}, y consultas punto-en-polígono cuyo resultado se verifica contra el municipio
de origen) se reportan en \texttt{docs/evidencia/spatial\_evidence.md}. <Comentar aquí: con 170 documentos la
diferencia de tiempo es pequeña, pero el plan cambia de \texttt{COLLSCAN} a \texttt{IXSCAN}; la ventaja crece en
las colecciones de edificios (Entrega 3), con millones de documentos.>

% (requiere booktabs y listings, ya cargados en main.tex). Reemplace los campos <...> con la salida real de
% src/spatial_evidence.py (docs/evidencia/spatial_evidence.md).
\section{Integración Espacial en MongoDB}

\subsection{Carga e índices}
Los municipios PDET limpios (GeoJSON en EPSG:4326, producido por el ETL de integridad) se cargan en la colección
\texttt{municipalities} con \texttt{src/load\_to\_mongo.py}. La carga es idempotente (\emph{upsert} por
\texttt{\_id}, el código DIVIPOLA), las geometrías se normalizan a \texttt{MultiPolygon} con la regla de la mano
derecha y se rechazan las inválidas o ubicadas fuera de la caja de Colombia. Los índices se crean antes de insertar,
de modo que MongoDB también rechaza cualquier geometría que no pueda indexar.

\begin{lstlisting}[style=lab,language=mongosh]
db.municipalities.createIndex({ geometry: "2dsphere" }, { name: "geometry_2dsphere" })
db.municipalities.createIndex({ pdet_subregion: 1 },    { name: "pdet_subregion_1" })
\end{lstlisting}

La carga queda registrada en \texttt{ingestion\_log} (archivo, filas leídas, insertadas y rechazadas, SHA-256 del
archivo y fecha). Resultado: <N> municipios cargados de <N> esperados (170), <S> subregiones (16), <R> rechazados.

\subsection{Evidencia del uso de los índices}
Se ejecutaron consultas \texttt{\$geoIntersects}, \texttt{\$geoWithin} y \texttt{\$near} con
\texttt{explain("executionStats")}, comparando el plan con el índice \texttt{2dsphere} frente a un recorrido
completo forzado con \texttt{hint(\{\$natural: 1\})}. El Cuadro~\ref{tab:evidencia} resume los planes ganadores y
los documentos examinados.

\begin{table}[!htbp]
\centering\small
\begin{tabular}{p{0.34\linewidth}p{0.28\linewidth}rrr}
\toprule
Consulta & Plan ganador & Docs (índice) & Docs (sin índice) & Devueltos \\
\midrule
\texttt{punto\_en\_poligono[05031]} & FETCH/IXSCAN & 1 & 170 & 1 \\
\texttt{punto\_en\_poligono[13688]} & FETCH/IXSCAN & 1 & 170 & 1 \\
\texttt{punto\_en\_poligono[19548]} & FETCH/IXSCAN & 2 & 170 & 1 \\
\texttt{punto\_en\_poligono[47001]} & FETCH/IXSCAN & 1 & 170 & 1 \\
\texttt{punto\_en\_poligono[54800]} & FETCH/IXSCAN & 1 & 170 & 1 \\
\texttt{punto\_en\_el\_oceano} & FETCH/IXSCAN & 0 & 170 & 0 \\
\texttt{ventana\_geoWithin} & FETCH/IXSCAN & 34 & 170 & 18 \\
\texttt{vecinos\_geoIntersects} & FETCH/IXSCAN & 5 & 170 & 1 \\
\texttt{cercanos\_near} & GEO\_NEAR\_2DSPHERE/FETCH/IXSCAN/FETCH/IXSCAN/FETCH/IXSCAN & 335 & -- & 18 \\
\bottomrule
\end{tabular}
   % generado por src/spatial_evidence.py
\caption{Planes de ejecución de las consultas espaciales sobre \texttt{municipalities}.}
\label{tab:evidencia}
\end{table}

Los controles automáticos (índice presente, número de municipios y subregiones, códigos DIVIPOLA de cinco
dígitos, tipo \texttt{MultiPolygon}, y consultas punto-en-polígono cuyo resultado se verifica contra el municipio
de origen) se reportan en \texttt{docs/evidencia/spatial\_evidence.md}. <Comentar aquí: con 170 documentos la
diferencia de tiempo es pequeña, pero el plan cambia de \texttt{COLLSCAN} a \texttt{IXSCAN}; la ventaja crece en
las colecciones de edificios (Entrega 3), con millones de documentos.>
